% TEMPLATE-MILC-v3.tex - plantilla maestra UNICA de las guias MILC v3.
%
% La usan las tres familias de guias, cada una con su driver:
%   · Grados 6-11  -> scripts/build-guias-g11.py
%   · Bebras       -> scripts/build-guias-bebras.py
%   · Semillero    -> scripts/build-guias-semillero.py
%
% Antes habia tres copias casi identicas (~400 lineas iguales, 6 distintas):
% cada mejora habia que aplicarla tres veces, y en la practica no se hacia
% (el motor de recursos vivia solo en dos de ellas). Lo que varia entre
% familias esta parametrizado en 6 placeholders que cada driver llena
% (se nombran aqui SIN su sintaxis real: el builder reemplaza por texto plano,
% tambien dentro de los comentarios, y un valor multilinea rompe el preambulo):
%
%   HEADER_IZQ        encabezado izquierdo de pagina
%   HEADER_DER        encabezado derecho de pagina
%   PORTADA_ETIQUETA  rotulo sobre el numero grande (GRADO/MOMENTO/MODULO)
%   PORTADA_SERIE     linea de serie de la portada
%   PORTADA_ID        identificador de la guia en la portada
%   PORTADA_CIRCULO   el circulito de grado (°), o vacio si no aplica
%
% El resto de claves las documenta content/guias/_SCHEMA.md. El script valida
% que no queden placeholders sin reemplazar antes de escribir el archivo final.

\documentclass[11pt,letterpaper]{article}
\usepackage[margin=1.75cm]{geometry}
\usepackage[table]{xcolor}
\usepackage{fontspec}
\usepackage[spanish]{babel}
\usepackage{tikz}
% Librerías TikZ para los diagramas de las guías (bloque `recursos:`):
%  · babel  → babel en español vuelve activos caracteres como «>», y TikZ los usa
%             en sus opciones (>=stealth, ->). Sin esta librería, cualquier
%             diagrama con flechas revienta con «\language@active@arg> has an
%             extra }». Debe ir ANTES de que se dibuje nada.
%  · positioning        → sintaxis «below left=2pt and 3pt».
%  · decorations.pathreplacing → llaves ({brace}) para señalar rangos.
\usetikzlibrary{babel,positioning,decorations.pathreplacing}
\usepackage{graphicx}
% Circuitos: el trabajo de electrónica se hace hoy con micro:bit y sus sensores
% integrados (MakeCode, sin cableado), así que no se necesitan esquemáticos.
% Si algún día entran montajes con protoboard, basta descomentar esta línea:
% los diagramas .tex/.tikz declarados en `recursos:` ya se incrustan solos.
% \usepackage{circuitikz}
\usepackage[most]{tcolorbox}
\usepackage{tabularx}
\usepackage{array}
\usepackage{fancyhdr}
\usepackage{titlesec}
\usepackage{ragged2e}
\usepackage{enumitem}
\usepackage{microtype}

% Helvetica Neue no trae la flecha U+2192, asi que cada «→» escrito literalmente
% en el YAML se compilaba a NADA, en silencio (462 flechas en 61 guias: rutas del
% tipo «paleta → tipografia → lockup» salian rotas). Mapeamos el caracter a su
% version matematica, que si tiene glifo. Asi el autor puede seguir escribiendo
% «→» comodamente en el YAML.
\usepackage{newunicodechar}
\newunicodechar{→}{\ensuremath{\rightarrow}}

\setmainfont{Helvetica Neue}
\setsansfont{Helvetica Neue}
\setlength{\parindent}{0pt}
\setlength{\parskip}{6pt}
\hyphenpenalty=200
\exhyphenpenalty=200
\pretolerance=400
\tolerance=2000
\setlength{\emergencystretch}{1em}
\renewcommand{\arraystretch}{1.25}
\setlist{leftmargin=5mm,itemsep=1mm,topsep=1mm}
\newcolumntype{Y}{>{\RaggedRight\arraybackslash}X}

\definecolor{milcMagenta}{HTML}{E6007E}
\definecolor{milcVino}{HTML}{5A0038}
\definecolor{milcTurquesa}{HTML}{0093A5}
\definecolor{milcVerde}{HTML}{58A923}
\definecolor{milcMostaza}{HTML}{E5B400}
\definecolor{milcOcre}{HTML}{B86600}
\definecolor{milcNegro}{HTML}{241816}
\definecolor{milcGris}{HTML}{F7F7F4}
\definecolor{milcLinea}{HTML}{E8E2DD}
\definecolor{milcGradePrimary}{HTML}{FF2D87}
\definecolor{milcGradeDark}{HTML}{9F1B56}
\definecolor{milcGradeSoft}{HTML}{FFE0EE}
\definecolor{milcGradeAccent}{HTML}{CFE7FF}
\definecolor{milcGradeText}{HTML}{FFFFFF}
\color{milcNegro}

\pagestyle{fancy}
\fancyhf{}
\lhead{\small Tecnología e Informática · Grado 10}
\rhead{\small Guía 11 · MILC v3}
\cfoot{\small\thepage}
\renewcommand{\headrulewidth}{0pt}

\newtcolorbox{titlebox}[1]{
  enhanced,colback=#1,colframe=#1,arc=7mm,boxrule=0pt,
  left=7mm,right=7mm,top=3mm,bottom=3mm,
  before skip=8pt,after skip=8pt,
  fontupper=\sffamily\bfseries\Large\color{white}
}

\newtcolorbox{softbox}[3]{
  enhanced,colback=#2,colframe=#1,borderline west={4pt}{0pt}{#1},
  arc=2mm,boxrule=.4pt,left=4mm,right=4mm,top=3mm,bottom=3mm,
  before skip=6pt,after skip=8pt,title={#3},coltitle=white,
  colbacktitle=#1,fonttitle=\sffamily\bfseries,fontupper=\RaggedRight,
  attach boxed title to top left={xshift=3mm,yshift=-2mm},
  boxed title style={arc=2mm, boxrule=0pt}
}

\newtcolorbox{drawbox}[2]{
  enhanced,colback=white,colframe=#1,arc=4mm,boxrule=1pt,
  left=5mm,right=5mm,top=4mm,bottom=4mm,before skip=8pt,after skip=8pt,
  title={#2},coltitle=white,colbacktitle=#1,fonttitle=\sffamily\bfseries,
  fontupper=\RaggedRight,
  attach boxed title to top left={xshift=4mm,yshift=-2mm},
  boxed title style={arc=3mm, boxrule=0pt}
}

\newtcolorbox{infoband}[1]{
  enhanced,colback=#1!8,colframe=#1!50,arc=2mm,boxrule=.5pt,
  left=4mm,right=4mm,top=2mm,bottom=2mm,before skip=4pt,after skip=4pt,
  fontupper=\small\RaggedRight
}

% Puente narrativo entre secciones (contrato editorial Conéctate).
% Da continuidad: "Acabas de X. Ahora vas a Y porque Z."
% Visualmente: banda izquierda en color del grado, texto en itálica pequeña.
\newtcolorbox{puentebox}{
  enhanced,colback=milcGradeSoft,colframe=milcGradeDark,
  borderline west={3pt}{0pt}{milcGradeDark},
  arc=0mm,boxrule=0pt,left=5mm,right=4mm,top=2.5mm,bottom=2.5mm,
  before skip=4pt,after skip=8pt,
  fontupper=\itshape\small\color{milcNegro}\RaggedRight
}
% La flecha va en modo matematico a proposito: Helvetica Neue NO tiene el glifo
% U+2192, asi que \textrightarrow se compilaba a NADA («Missing character: There
% is no → in font Helvetica Neue») y el puente salia sin flecha. En modo
% matematico la flecha viene de la fuente matematica y si se dibuja.
\newcommand{\puente}[1]{%
  \begin{puentebox}%
    \textcolor{milcGradeDark}{$\rightarrow$}\hspace{2mm}#1%
  \end{puentebox}%
}

\newcommand{\linead}[1]{\noindent\color{milcLinea}\rule{#1}{0.5pt}\color{milcNegro}}
\newcommand{\respuesta}[1]{\par\vspace{2mm}\foreach \n in {1,...,#1}{\noindent\color{milcLinea}\rule{\linewidth}{0.5pt}\par\vspace{5mm}}\color{milcNegro}}
\newcommand{\checkbox}{\textcolor{milcGradeDark}{\fbox{\phantom{x}}}\hspace{5pt}}

\newcommand{\phasecard}[4]{
\begin{tcolorbox}[enhanced,colback=#3,colframe=#2,arc=3mm,boxrule=.5pt,height=3.05cm,valign=center,halign=center,left=1.5mm,right=1.5mm,top=2mm,bottom=2mm]
{\sffamily\bfseries\fontsize{22}{24}\selectfont\textcolor{#2}{#1}}\par
{\sffamily\bfseries\small #4}
\end{tcolorbox}
}

% ── Recursos visuales de la guía ──────────────────────────────────────
% Declarados en el bloque `recursos:` del YAML y emitidos por el builder.
% \guiaFigura   → imagen rasterizada o PDF (foto, carta celeste, montaje).
% \guiaDiagrama → fuente TikZ/circuitikz (.tex) incrustada como vector:
%                 es la vía para circuitos y esquemáticos editables.
% #1 = ruta relativa al .tex · #2 = pie (puede ir vacío)
\newcommand{\guiaFigura}[2]{%
\begin{center}
\includegraphics[width=0.88\linewidth,height=0.38\textheight,keepaspectratio]{#1}\\[1.5mm]
{\footnotesize\itshape\textcolor{milcNegro!75}{#2}}
\end{center}
}

\newcommand{\guiaDiagrama}[2]{%
\begin{center}
\input{#1}\\[1.5mm]
{\footnotesize\itshape\textcolor{milcNegro!75}{#2}}
\end{center}
}

\begin{document}
\thispagestyle{empty}
\begin{tikzpicture}[remember picture,overlay]
  \fill[white] (current page.south west) rectangle (current page.north east);
  \path[fill=milcGradePrimary,rounded corners=16mm]
    ([xshift=.55cm,yshift=-.55cm]current page.north west) rectangle
    ([xshift=-.55cm,yshift=3.65cm]current page.south east);

  \node[anchor=north west,text=milcGradeText] at ([xshift=2.0cm,yshift=-1.85cm]current page.north west)
    {\sffamily\bfseries\fontsize{14}{16}\selectfont GRADO};
  \node[anchor=north west,text=milcGradeText,opacity=.82] at ([xshift=2.0cm,yshift=-2.75cm]current page.north west)
    {\sffamily\bfseries\fontsize{9}{11}\selectfont SERIE GUÍAS · TECNOLOGÍA E INFORMÁTICA};

  \begin{scope}[shift={([xshift=-3.05cm,yshift=-2.55cm]current page.north east)}]
    \fill[white,opacity=.80] (-.62,-.52) rectangle (.62,.52);
    \draw[milcGradeText,line width=1pt] (-.62,-.52) rectangle (.62,.52);
    \fill[milcGradeDark] (-.42,-.38) rectangle (-.22,.06);
    \fill[milcGradeAccent] (-.10,-.38) rectangle (.10,.24);
    \fill[milcGradePrimary!60!black] (.22,-.38) rectangle (.42,.38);
  \end{scope}

  \node[anchor=west,text=milcGradeText] at ([xshift=1.9cm,yshift=-12.85cm]current page.north west)
    {\sffamily\bfseries\fontsize{124}{124}\selectfont 10};
  % El circulito de grado (°) solo aplica a las guias de grado («11°»); en
  % Bebras y Semillero el numero grande es un momento/modulo, no un grado.
  \draw[milcGradeText,line width=1.7pt]
    ([xshift=4.52cm,yshift=-11.68cm]current page.north west) circle (.13cm);

  \node[anchor=west,text=milcGradeText,text width=8.7cm,align=left] at ([xshift=10.35cm,yshift=-11.6cm]current page.north west)
    {
     {\sffamily\bfseries\fontsize{19}{22}\selectfont 4 tipos de texto\\profesional ---\\cada oficio su tono}\\[4mm]
     {\sffamily\bfseries\fontsize{23}{26}\selectfont Guía 11-10-TIC}\\[2mm]
     {\sffamily\fontsize{13}{16}\selectfont Período 2 · Informes técnicos y comunicación profesional con IA}\\[5mm]
     {\sffamily\bfseries\fontsize{11}{13}\selectfont Institución Educativa Sor María Juliana}\\[1mm]
     {\sffamily\fontsize{10}{12}\selectfont Autor: Dr. Álvaro Cárdenas Orozco}
    };

  \path[fill=milcGradeSoft,rounded corners=7mm]
    ([xshift=1.35cm,yshift=.85cm]current page.south west) rectangle
    ([xshift=-1.35cm,yshift=4.25cm]current page.south east);
  \draw[milcGradeDark,line width=.65pt,rounded corners=7mm]
    ([xshift=1.35cm,yshift=.85cm]current page.south west) rectangle
    ([xshift=-1.35cm,yshift=4.25cm]current page.south east);
  \node[anchor=center,text=milcNegro,text width=17.0cm,align=center] at ([yshift=2.55cm]current page.south)
    {\sffamily\fontsize{10}{12}\selectfont Metodología MILC · Apertura ancestral · Pensamiento computacional · Triángulo Dussel-Estoicismo-Floridi\\
    \textcolor{milcGradeDark}{\bfseries Producto:} Tabla comparativa de los 4 tipos de texto profesional con (a) propósito específico, (b) audiencia típica, (c) estructura mínima, (d) tono apropiado + 1 ejemplo real por cada tipo capturado de prensa, redes profesionales (LinkedIn), publicaciones institucionales o publicidad + análisis breve de 5 líneas explicando qué diferencia concretamente un informe técnico de uno comercial + reflexión sobre cuál de los 4 tipos vas a usar más en tu vida laboral próxima};
\end{tikzpicture}
\mbox{}
\newpage

\section*{Apertura: 4 tipos de texto profesional --- informe, comercial, técnico, estudio de mercado}

\begin{softbox}{milcGradeDark}{milcGradeSoft}{Saber ancestral}
En el barrio Obrero de Cartago hubo \textbf{cuatro oficios de la palabra} que cualquier abuela distinguía al oído: (1) \textbf{El cuentista}: el que contaba historias a los nietos. Su palabra era \emph{para entretener y enseñar con ejemplo}. Tono cálido, ritmo variado, metáforas locales. (2) \textbf{El escribano}: el señor con sombrero que redactaba cartas y poderes frente a la notaría. Su palabra era \emph{para acreditar formalmente}. Tono respetuoso, frases largas, fórmulas legales precisas. (3) \textbf{El contador}: el que llevaba las cuentas del tendero, el zapatero, la modista. Su palabra era \emph{para calcular y reportar con números}. Tono seco, cifras exactas, sin adornos. (4) \textbf{El juez}: el que oficiaba en la inspección de policía cuando había pleito vecinal. Su palabra era \emph{para sentenciar con autoridad}. Tono firme, argumentación clara. Cualquier vecino mayor distinguía los 4 al instante: \textbf{cada oficio tenía tono propio}. Confundirlos era perder el oficio. Los 4 tipos de texto profesional moderno son la versión actualizada de esos 4 oficios del barrio.
\end{softbox}

\begin{softbox}{milcTurquesa}{milcGris}{Saber contemporáneo}
En el mundo profesional contemporáneo, hay \textbf{4 tipos de texto} que cualquier persona con vida laboral encontrará una y otra vez: (1) \textbf{Informe técnico}: documento que reporta datos, análisis, hallazgos. Propósito: informar a quien necesita tomar decisiones técnicas. Audiencia: especialistas, supervisores, comités. Estructura: introducción + metodología + resultados + conclusiones + recomendaciones. Tono: neutro, datos verificables, sin adornos. (2) \textbf{Texto comercial}: propuesta de venta, presentación de servicio, pitch. Propósito: persuadir a un cliente o inversionista a comprar, contratar o invertir. Audiencia: cliente potencial. Estructura: problema + solución + beneficios + propuesta + llamada a la acción. Tono: cálido pero profesional, claro, orientado al cierre. (3) \textbf{Texto comercial-técnico mixto} o estudio: combinación que se ve en propuestas a empresas con componente técnico. (4) \textbf{Estudio de mercado}: investigación que analiza un sector, audiencia o competencia. Propósito: dar contexto para decisiones estratégicas. Audiencia: gerencia, emprendedores, inversionistas. Estructura: contexto + metodología + análisis + conclusiones. Tono: analítico, basado en datos, neutro. La regla profesional: \textbf{usar el tono del oficio equivocado pierde la audiencia}. Un informe técnico con tono comercial parece superficial; un texto comercial con tono técnico no vende.
\end{softbox}

\begin{softbox}{milcMostaza}{milcGris}{Pregunta-puente}
\textit{¿Qué sabían las abuelas del barrio al distinguir al cuentista del escribano y del contador, que el novato olvida cuando escribe todo con el mismo tono? ¿Y por qué un mismo dato presentado en informe técnico vs en texto comercial requiere palabras y estructuras completamente distintas?}
\end{softbox}

\begin{softbox}{milcVerde}{milcGris}{Saber y saber hacer}
Al terminar podrás: (1) \textbf{identificar} cuál de los 4 tipos de texto profesional aplica a una situación concreta del mundo real; (2) \textbf{analizar} ejemplos reales detectando qué tono y estructura sostiene cada tipo; (3) \textbf{explicar} con tus palabras la diferencia entre informe técnico y comercial usando un mismo tema como ejemplo; (4) \textbf{evaluar} cuál de los 4 vas a usar más en tu próxima vida laboral y por qué.
\end{softbox}



\small
\begin{tabularx}{\textwidth}{>{\bfseries}p{3.0cm}Y}
\rowcolor{milcGradePrimary!12}Autor & Dr. Álvaro Cárdenas Orozco\\
DBA & Comunicación profesional con asistencia de IA y herramientas ofimáticas gratuitas (MEN, T\&I)\\
Referentes & Markdown como estándar abierto · Google Workspace · Estoicismo (claridad)\\
Producto final & Tabla comparativa de los 4 tipos de texto profesional con (a) propósito específico, (b) audiencia típica, (c) estructura mínima, (d) tono apropiado + 1 ejemplo real por cada tipo capturado de prensa, redes profesionales (LinkedIn), publicaciones institucionales o publicidad + análisis breve de 5 líneas explicando qué diferencia concretamente un informe técnico de uno comercial + reflexión sobre cuál de los 4 tipos vas a usar más en tu vida laboral próxima\\
\end{tabularx}
\normalsize

\puente{Antes de empezar, mira el viaje completo. Cerraste el proyecto del libro (P1). Hoy abre el periodo 2: \textbf{texto profesional con asistencia de IA}. Las próximas 9 sesiones (estructura del informe, Google Docs, Markdown, encuestas, análisis con IA, informe comercial, correo profesional, mini-proyecto, sustentación) construyen sobre la habilidad de distinguir los 4 tipos.}

\begin{titlebox}{milcGradeDark}
Ruta MILC: así vamos a aprender
\end{titlebox}
\begin{tabularx}{\textwidth}{>{\centering\arraybackslash}X>{\centering\arraybackslash}X>{\centering\arraybackslash}X>{\centering\arraybackslash}X}
\phasecard{1}{milcVerde}{milcGris}{Escucha\\[-1mm]\small Observo, escucho y pregunto.} &
\phasecard{2}{milcTurquesa}{milcGris}{Sistematización\\[-1mm]\small Pensamiento computacional.} &
\phasecard{3}{milcMagenta}{milcGris}{Praxis\\[-1mm]\small Construyo y aplico.} &
\phasecard{4}{milcVino}{milcGris}{Evaluación\\[-1mm]\small Reflexión liberadora.}
\end{tabularx}


\puente{Antes de teorizar, vas a hacer una \emph{auditoría rápida}: busca 4 textos reales (en LinkedIn, prensa, sitios institucionales, publicidad). Para cada uno, intenta identificar qué tipo es. Esa observación intuitiva te enseña los 4 tonos en uso.}

\begin{titlebox}{milcVerde}
Fase 1 · Escucha
\end{titlebox}

\begin{softbox}{milcVerde}{milcGris}{Escena para escuchar}
\textbf{Actividad 1 · IDENTIFICA --- Auditoría de 4 textos reales} (15 min · individual). Busca 4 textos reales en internet, redes profesionales (LinkedIn), prensa o publicidad. Idealmente uno de cada tipo: (1) un \textbf{informe técnico} (informe del DANE, informe de gestión escolar, reporte de Cámara de Comercio). (2) un \textbf{texto comercial} (propuesta comercial publicada, presentación de servicio, anuncio publicitario detallado). (3) un \textbf{estudio de mercado} (análisis sectorial, reporte de tendencias publicado). (4) un \textbf{informe técnico-comercial mixto} (caso de éxito empresarial, white paper). Para cada uno anota: ¿qué tipo crees que es?, ¿qué te indica el tono?.
\end{softbox}

\checkbox Encontré 4 textos reales (1 por tipo).\\
\checkbox Identifiqué qué tipo es cada uno.\\
\checkbox Reconocí qué tono distingue a cada tipo.

\begin{infoband}{milcVerde}


\textbf{Lo que va al cuaderno (Actividad 1 · IDENTIFICA).} Título: «Actividad 1 --- 4 textos reales». \textbf{Formato:} 4 fichas (1 por texto) con link/captura + tipo identificado + tono observado. \textbf{Sabes que terminaste cuando} distingues los 4 al primer párrafo.
\end{infoband}

\puente{Ya tienes 4 ejemplos. Ahora vas a entender la \textbf{anatomía de los 4 tipos}: propósito, audiencia, estructura, tono. Y los 5 errores típicos de confundir tipos.}

\begin{titlebox}{milcTurquesa}
Fase 2 · Sistematización con pensamiento computacional
\end{titlebox}

\begin{softbox}{milcTurquesa}{milcGris}{Cuatro pilares aplicados al tema}
Tus 4 ejemplos te dan vocabulario. Ahora necesitas la \textbf{anatomía profesional} de los 4 tipos.
\end{softbox}

\small
\begin{tabularx}{\textwidth}{>{\bfseries\RaggedRight\arraybackslash}p{3.6cm}Y}
\rowcolor{milcTurquesa!90}\color{white}Pilar & \color{white}\bfseries Aplicación al tema\\
1. Descomposición & Los \textbf{4 tipos de texto profesional} se descomponen así: (1) \textbf{Informe técnico}: \emph{propósito} informar para decisión técnica; \emph{audiencia} especialistas o supervisores; \emph{estructura} intro + metodología + resultados + conclusiones + recomendaciones; \emph{tono} neutro, datos verificables, sin opiniones personales. Ejemplos cotidianos: informe de gestión escolar, reporte de proyecto. (2) \textbf{Texto comercial}: \emph{propósito} persuadir para vender o contratar; \emph{audiencia} cliente potencial; \emph{estructura} problema + solución + beneficios + propuesta + llamada a la acción; \emph{tono} cálido pero profesional, orientado al cierre. (3) \textbf{Estudio de mercado}: \emph{propósito} dar contexto estratégico; \emph{audiencia} gerencia, emprendedores, inversionistas; \emph{estructura} contexto sectorial + metodología + análisis + conclusiones; \emph{tono} analítico, basado en datos. (4) \textbf{Texto técnico-comercial mixto}: \emph{propósito} explicar técnicamente con propósito comercial; \emph{audiencia} cliente sofisticado; \emph{estructura} problema técnico + solución detallada + caso de éxito + propuesta; \emph{tono} riguroso pero accesible.\\
2. Reconocimiento de patrones & El patrón profesional clave: \textbf{el tono se elige por la audiencia, no por preferencia personal}. Si tu lector es un ingeniero, el tono técnico le inspira confianza. Si tu lector es un cliente novato, el tono técnico lo aleja. La phronesis del oficio consiste en \textbf{ajustar el registro a quien lee, no al que escribe}.\\
3. Abstracción & La \textbf{diferencia clave entre informe técnico y comercial} con el mismo tema: ej. \emph{``nueva máquina de café''}. (a) Informe técnico: \emph{``La cafetera modelo X-200 procesa 50 tazas/hora con consumo 1.2 kWh, presión 9 bar, tanque 3L. Tiempo medio entre fallos según pruebas: 1200 horas''}. (b) Texto comercial: \emph{``Imagine empezar cada mañana con café de cafetería en su casa. La X-200 lo hace posible en 60 segundos. Diseño elegante, manejo simple, garantía 2 años. Hoy con 20 por ciento de descuento''}. Mismo producto, distinta audiencia, distinto tono, distinta estructura.\\
4. Algoritmo conceptual & Algoritmo de identificación: (1) leer el primer párrafo. (2) preguntar: ¿quién es el lector ideal? (3) preguntar: ¿qué quiere lograr el texto? (4) según la combinación lector + propósito, identificar el tipo. Una vez identificado, ajusta tu lectura: en informe técnico busca datos; en texto comercial busca beneficios; en estudio busca tendencias; en mixto busca rigor con cierre.\\
\end{tabularx}
\normalsize

\begin{drawbox}{milcTurquesa}{Anatomía de los 4 tipos}
\textbf{Informe técnico:} neutro, datos, decisión técnica. \\[1mm] \textbf{Comercial:} cálido, persuasivo, cierre de venta. \\[1mm] \textbf{Estudio de mercado:} analítico, contexto estratégico. \\[1mm] \textbf{Técnico-comercial mixto:} riguroso + accesible.
\end{drawbox}

\begin{softbox}{milcMostaza}{milcGris}{Errores comunes y su corrección}
(1) \textbf{Tono comercial en informe técnico}: el informe parece publicidad y pierde credibilidad. (2) \textbf{Tono técnico en texto comercial}: el cliente novato no entiende y no compra. (3) \textbf{Sin estructura clara}: mezcla de todo en orden caprichoso. (4) \textbf{Audiencia mal identificada}: escribir para especialista cuando lee novato. (5) \textbf{Conclusiones sin datos en informe}: opinión disfrazada de análisis. (6) \textbf{Beneficios sin emoción en texto comercial}: lista seca de características sin conexión.
\end{softbox}

\begin{infoband}{milcTurquesa}


\textbf{Lo que va al cuaderno (Actividad 2 · ANALIZA).} Título: «Actividad 2 --- Anatomía 4 tipos». \textbf{Formato:} (a) los 4 tipos con propósito/audiencia/estructura/tono; (b) ejemplo informe vs comercial sobre mismo tema; (c) los 6 errores con marca. \textbf{Sabes que terminaste cuando} reconoces el tipo al primer párrafo.
\end{infoband}

\puente{Tienes el método. Toca aplicarlo: construyes tabla comparativa, agregas 1 ejemplo real por tipo, escribes análisis de diferencia informe vs comercial, reflexionas sobre cuál usarás más.}

\begin{titlebox}{milcMagenta}
Fase 3 · Praxis con pensamiento computacional
\end{titlebox}

\begin{softbox}{milcMagenta}{milcGris}{Construyendo el producto con los cuatro pilares}
\textbf{Actividad 3 · EXPLICA + EVALÚA --- Tabla + ejemplos + análisis + reflexión} (35 min · individual). Hora de aplicar. Construyes tabla comparativa, agregas 1 ejemplo real por cada tipo, escribes análisis de diferencia, reflexionas sobre cuál usarás más en tu vida laboral.
\end{softbox}

\small
\begin{tabularx}{\textwidth}{>{\bfseries\RaggedRight\arraybackslash}p{3.6cm}Y}
\rowcolor{milcMagenta!90}\color{white}Pilar & \color{white}\bfseries Aplicación al producto\\
Descomposición & Tu tabla comparativa tiene 4 filas (1 por tipo) × 5 columnas (Tipo, Propósito, Audiencia, Estructura, Tono). Llena cada celda con 1-2 frases concretas. Para cada tipo, agrega también un link o referencia al ejemplo real que encontraste en la Actividad 1.\\
Patrones & El análisis de diferencia informe vs comercial se escribe en 5 líneas usando un mismo tema (puede ser un producto, un servicio, un evento). Muestra cómo cambia la palabra cuando cambia la audiencia. Ejemplo: \emph{``Una computadora portátil. En informe técnico: marca, modelo, especificaciones técnicas, consumo, garantía técnica. En texto comercial: \emph{lleva tu trabajo a cualquier parte}, batería todo el día, diseño elegante, precio promocional''}.\\
Abstracción & La reflexión final responde: \emph{¿cuál de los 4 tipos voy a usar más en mi vida laboral próxima?}. Esto depende de tu vocación: si vas a estudiar ingeniería, probablemente informes técnicos. Si emprendimiento, textos comerciales. Si administración pública, informes y estudios. La reflexión te ayuda a priorizar qué texto profesional dominar mejor en las próximas sesiones.\\
Algoritmo de armado & Algoritmo: (1) construir tabla 4 filas × 5 columnas. (2) agregar 1 ejemplo real por tipo. (3) escribir análisis de 5 líneas comparando informe vs comercial sobre mismo tema. (4) escribir reflexión sobre qué tipo usarás más. (5) verificar coherencia entre tabla y análisis.\\
\end{tabularx}
\normalsize

\begin{drawbox}{milcMagenta}{Checklist antes de entregar}
\checkbox Tabla con 4 tipos y 5 columnas. \\[1mm] \checkbox 1 ejemplo real por cada tipo (link o captura). \\[1mm] \checkbox Análisis informe vs comercial con mismo tema. \\[1mm] \checkbox Reflexión sobre qué tipo usarás más. \\[1mm] \checkbox Tabla coherente con análisis. \\[1mm] \checkbox Distingo los 4 tipos al primer párrafo.
\end{drawbox}

\begin{softbox}{milcMagenta}{milcGris}{Plantilla de trabajo}
\textbf{Informe técnico.} \textit{Propósito + audiencia + estructura + tono + ejemplo.} \\[1mm] \textbf{Texto comercial.} \textit{Idem.} \\[1mm] \textbf{Estudio de mercado.} \textit{Idem.} \\[1mm] \textbf{Mixto técnico-comercial.} \textit{Idem.} \\[1mm] \textbf{Diferencia informe vs comercial.} \textit{Mismo tema, distinta palabra.} \\[1mm] \textbf{Reflexión.} \textit{Cuál usaré más y por qué.}
\end{softbox}

\begin{infoband}{milcMagenta}


\textbf{Lo que va al cuaderno (Actividad 3 · EXPLICA + EVALÚA).} Título: «Actividad 3 --- Los 4 oficios de la palabra». \textbf{Formato:} tabla + 4 ejemplos + análisis + reflexión. \textbf{Extensión:} 2-3 páginas. \textbf{Sabes que terminaste cuando} podrías identificar al instante qué tipo te piden cuando te pidan \emph{``escríbeme un texto sobre X''}.
\end{infoband}

\puente{Lo que sigue resume \emph{qué} entregas hoy: tabla + 4 ejemplos + análisis + reflexión. El criterio principal: que puedas reconocer al instante qué tipo de texto te están pidiendo cuando alguien te diga \emph{``escríbeme un X''}.}

\begin{titlebox}{milcMostaza}
Producto de la sesión
\end{titlebox}
\textbf{Tabla 4 tipos + 4 ejemplos + análisis + reflexión personal}.
\begin{softbox}{milcMostaza}{milcGris}{Criterios de calidad}
(1) Tabla comparativa con 5 columnas por tipo. (2) 1 ejemplo real por cada tipo. (3) Análisis informe vs comercial con mismo tema. (4) Reflexión sobre uso futuro. (5) Coherencia tabla con análisis. (6) Distinción de tipos al primer párrafo.
\end{softbox}

\puente{Antes de cerrar, mira los 4 tipos desde las cinco dimensiones humanas. Distinguir los tonos es respeto por la audiencia de cada texto; usar uno solo para todo es perder oficio profesional.}

\begin{titlebox}{milcVino}
Fase 4 · Evaluación liberadora — cinco dimensiones
\end{titlebox}

\begin{softbox}{milcVino}{milcGris}{Sentido de la evaluación}
Evaluar no es solo revisar una nota. Es reconocer qué aprendí, qué cambió en mí, qué emociones manejé, cómo me relaciono con la ciudad y la comunidad, y qué le dejaré a quien viene detrás.
\end{softbox}

\textbf{Mi reflexión liberadora en cinco dimensiones.} Completa cada frase iniciada:

\small
\begin{tabularx}{\textwidth}{>{\bfseries\RaggedRight\arraybackslash}p{3.4cm}Y>{\RaggedRight\arraybackslash}p{4.6cm}}
\rowcolor{milcVino}\color{white}Dimensión & \color{white}\bfseries Mi respuesta (frame starter) & \color{white}\bfseries Algo que descubrí\\
1. Personal & Lo que más cambió en mí fue \linead{6.5cm} & \linead{4.2cm}\\
2. Emocional & La emoción más fuerte fue \linead{2.5cm} y la regulé \linead{3cm} & \linead{4.2cm}\\
3. Ciudadana & En democracia esto importa porque \linead{6cm} & \linead{4.2cm}\\
4. Local (Cartago) & En mi barrio sería útil para \linead{6.5cm} & \linead{4.2cm}\\
5. Intergeneracional & A mi abuela/abuelo le diría \linead{6.5cm} & \linead{4.2cm}\\
\end{tabularx}
\normalsize

\begin{infoband}{milcVino}
\textbf{Para la próxima sustentación.} Vuelve a esta tabla en seis meses. Verás que las respuestas cambian — no porque lo aprendido cambie, sino porque tú cambias. La evaluación liberadora es una conversación contigo a través del tiempo.
\end{infoband}


\puente{Y para terminar, tres voces sobre la palabra adecuada. Dussel sobre los textos que sirven a audiencias específicas, Séneca sobre la sobriedad del tono correcto, Floridi sobre la comunicación contextual como ética profesional.}

\begin{titlebox}{milcGradeDark}
Triángulo de pensamiento · cierre filosófico
\end{titlebox}

\begin{softbox}{milcVino}{milcGris}{Dussel · lente del nosotros}
\textit{``El tono que respeta a la audiencia es liberador; el tono uniforme impuesto a toda audiencia es violencia comunicativa.''}

Cuando escribes con el mismo tono para todas las audiencias, estás imponiendo tu manera de hablar a quien debe leer. La filosofía de la liberación pide lo opuesto: \textbf{ajustar el registro a quien lee}. El cliente novato merece tono accesible; el especialista merece tono técnico. La phronesis comunicativa consiste en \textbf{servir a la audiencia, no imponer la propia voz}.

\textbf{¿Mi escritura ajusta el tono a quien lee, o uso el mismo registro para todo?}
\end{softbox}
\linead{\linewidth}\\[1mm]
\linead{\linewidth}

\begin{softbox}{milcMostaza}{milcGris}{Estoicismo · Séneca · cuidado interior}
\textit{``La sobriedad del tono correcto es virtud profesional; el adorno inadecuado a la audiencia es vicio del que escribe.''}

Es tentador querer sonar elegante en todo texto. La sabiduría estoica recomienda lo opuesto: \emph{usa el tono justo para la audiencia, ni más ni menos}. Un informe técnico con prosa florida pierde credibilidad. Un texto comercial con frialdad técnica pierde clientes. La phronesis del oficio consiste en \textbf{aceptar la disciplina del registro apropiado}.

\textbf{¿Estoy usando el tono sobrio que cada tipo de texto exige, o lo adorno cuando no toca?}
\end{softbox}
\linead{\linewidth}\\[1mm]
\linead{\linewidth}

\begin{softbox}{milcTurquesa}{milcGris}{Floridi · lente de la infoesfera}
\textit{``La comunicación contextual es la nueva ética profesional en la era de la información distribuida.''}

En la era de la información, los textos circulan en contextos múltiples. La ética profesional pide que cada texto se ajuste al contexto de su lectura. Tu tabla comparativa de hoy te entrena en esa práctica que rige toda profesión que produzca documentos: \textbf{distinguir el tipo de texto según el contexto y la audiencia}.

\textbf{¿Sé reconocer en qué contexto y para quién escribe cada texto profesional que produzco?}
\end{softbox}
\linead{\linewidth}\\[1mm]
\linead{\linewidth}

\begin{titlebox}{milcVino}
Compromiso final
\end{titlebox}

\small
\begin{tabularx}{\textwidth}{>{\RaggedRight\arraybackslash}p{3.0cm}YYY}
\rowcolor{milcVino}\color{white}\bfseries Criterio & \color{white}\bfseries Lo logré & \color{white}\bfseries En proceso & \color{white}\bfseries Necesito apoyo\\
Comprensión del tema & Explico ideas centrales con mis palabras. & Reconozco ideas pero falta precisión. & Necesito repasar conceptos base.\\
Pensamiento computacional & Apliqué los 4 pilares al diseñar mi producto. & Apliqué algunos pilares. & Necesito releer la fase 2.\\
Aplicación y producto & Mi producto cumple los criterios y se puede revisar. & Producto funcional con ajustes pendientes. & Aún en construcción.\\
Reflexión 5 dimensiones & Respondí cada dimensión con honestidad. & Respondí algunas con detalle. & Mis respuestas son muy generales.\\
Triángulo de pensamiento & Conecté las 3 voces con mi trabajo real. & Conecté al menos una. & No relaciono las voces con mi proceso.\\
\end{tabularx}
\normalsize

\textbf{Hoy entendí que:}\\[1mm]
\linead{\linewidth}\\[1mm]
\linead{\linewidth}

\textbf{Mi compromiso para la próxima sesión es:}\\[1mm]
\linead{\linewidth}\\[1mm]
\linead{\linewidth}

\begin{softbox}{milcMostaza}{milcGris}{Cita de despedida · Marco Aurelio}
\textit{``El obstáculo en el camino se vuelve el camino. Lo que estorba al camino se vuelve camino.''}\\[1mm]
Cada error de hoy, bien observado, se convierte en pieza del camino que pisarás mañana.
\end{softbox}

\small
\begin{tabularx}{\textwidth}{>{\bfseries}p{3.6cm}Y}
\rowcolor{milcGradeDark!12}Nombre completo: & \linead{10cm}\\
Fecha: & \linead{4cm} \quad Curso/grupo: \linead{4cm}\\
Mi firma: & \linead{9cm}\\
\end{tabularx}
\normalsize

\begin{infoband}{milcGradeDark}
\textbf{Para la próxima sesión.} Llega con tu tabla y los 4 ejemplos. En la sesión 2 vas a aprender la \textbf{estructura del informe técnico}: introducción + desarrollo + conclusiones, aplicada a un caso real del colegio. Los tipos de hoy son el mapa; el informe técnico de mañana es la primera implementación específica.
\end{infoband}

\end{document}
